\chapter{Auditing Compatibility with xna4-spec}
\label{ch:xna4-spec-auditing}

Every coverage percentage this book has cited --- 227 of 245 public FNA types, 92.7\%,
Chapter~\ref{ch:what-is-cna} --- came from comparing CNA against a ground truth rather than
against memory or against FNA's own source. This chapter covers that ground truth directly:
\texttt{xna4-spec}, and what auditing against it actually looks like.

\section{What xna4-spec actually is}

\texttt{xna4-spec} is a complete, machine-readable XML transcription of the official XNA~4.0
documentation, originally published at Microsoft's own
\texttt{learn.microsoft.com/.../previous-versions/windows/xna/} archive --- every class,
struct, enum, interface, and delegate across all nineteen XNA namespaces, 544 types in total,
with each type's original property, method, constructor, field, and event descriptions
preserved. A master XSD schema and a top-level index of all 544 types organize the whole
collection; one directory exists per namespace, mirroring the real
\texttt{Microsoft.Xna.Framework.*} namespace tree exactly.

Per-namespace type counts give a sense of where the real API surface concentrates:
\cnans{Microsoft.Xna.Framework.Graphics} alone accounts for 175 of the 544 types, with
\cnans{GamerServices} (53) and \cnans{Content.Pipeline.Graphics} (46) the next largest ---
which is exactly why this book devoted an entire Part to Graphics and this Part gives
GamerServices its own chapter, rather than the namespace breakdown being an arbitrary
organizing choice on this book's part.

\section{Why this matters more than diffing against FNA}

Chapter~\ref{ch:media-system} already gave the sharpest illustration of why this distinction
is not academic: eight separate audit rounds of CNA's \cnaclass{Song} class diffed against
FNA's own source and missed a real gap (FNA itself omits \texttt{Album}, \texttt{Artist},
\texttt{Genre}, and \texttt{ToString()} relative to the real reference assemblies), because
every round trusted FNA as a stand-in for the real specification rather than checking the
real specification directly. \texttt{xna4-spec} exists precisely to make that class of
mistake avoidable: a diff against it checks CNA against the actual, original Microsoft
documentation, with FNA's own gaps and quirks no longer able to hide a CNA gap behind them.

\section{Running an audit yourself}

The practical shape of an audit, drawn from how CNA's own coverage figures are produced
(Chapter~\ref{ch:project-practice}'s \texttt{AUDIT.md} conventions), is a straightforward
three-step comparison: pick a namespace or class from \texttt{xna4-spec}'s own index, read
its full member list (properties, methods, constructors, fields, events) from the
corresponding XML file, and check each member off against CNA's actual header for that class
--- present and correctly shaped, present but tagged \texttt{NOXNA} (and therefore not really
part of what you are checking), or missing. A member present in CNA but absent from
\texttt{xna4-spec} and not tagged \texttt{NOXNA} at all is the specific ``Extra-unmarked''
finding category Chapter~\ref{ch:devices-sensors} showed a real instance of
(\cnaclass{SensorBase}'s \texttt{TimeBetweenUpdatesChanged} event, found and then correctly
retagged) --- exactly the class of bug the \texttt{NOXNA} compile-time gate from
Chapter~\ref{ch:testing-philosophy} exists to catch mechanically, with a manual
\texttt{xna4-spec} audit as the complementary, judgment-requiring check for anything that
gate cannot see on its own (a member that is tagged correctly but has the wrong signature, for
instance).

\section{A licensing note worth stating plainly}

\texttt{xna4-spec}'s own README is direct about its content's origin: the API names,
descriptions, and everything else in the collection are derived from Microsoft's own XNA~4.0
documentation and remain Microsoft's intellectual property. The repository presents itself as
an unofficial, community-maintained conversion into a machine-readable format, intended
solely for non-commercial, informational, and educational use --- the same spirit in which
this book cites specific facts drawn from it.

\section{What an audit cannot tell you}

One limit is worth naming directly: \texttt{xna4-spec} tells you whether a member exists and
what it is documented to do --- it cannot tell you whether CNA's own implementation of a
present, correctly-shaped member actually behaves correctly. That is exactly the job the
pixel tests, differential testing, and oracle methodology from
Chapter~\ref{ch:testing-philosophy} exist to do instead. A useful mental model for the rest of
this Part: \texttt{xna4-spec} auditing answers ``is the surface complete and correctly
shaped,'' while the testing mechanisms from the previous chapter answer ``does the surface
that exists actually work.'' Both questions matter, and neither answers the other.
